% fonts.tex — polices du PDF, choisies selon la machine.
%
% Les polices voulues sont celles de Microsoft : Calibri, Arial, Courier New.
% Elles ne sont pas installées partout — sous Linux, elles manquent le plus
% souvent —, et LuaLaTeX s'arrête sur une police introuvable plutôt que d'en
% prendre une autre. Chaque police est donc cherchée dans l'ordre, et la
% première présente est retenue :
%
%   texte     Calibri      -> Carlito          -> Liberation Sans
%   sans      Arial        -> Liberation Sans  -> DejaVu Sans
%   chasse    Courier New  -> Liberation Mono  -> DejaVu Sans Mono
%
% Carlito et les Liberation ont les mêmes dimensions que les polices qu'elles
% remplacent : la mise en page — coupures de lignes, de pages — ne change pas
% d'une machine à l'autre. Paquets Debian/Ubuntu : fonts-crosextra-carlito,
% fonts-liberation.
%
% Ces trois choix remplacent mainfont, sansfont et monofont des fichiers
% defaults-*.yaml : ils ne doivent pas y revenir, sans quoi Pandoc exigerait de
% nouveau la police nommée, présente ou non.
%
% Inclus par defaults-single.yaml et defaults-book.yaml (include-in-header),
% donc après le chargement de fontspec par le modèle LaTeX de Pandoc.

\IfFontExistsTF{Calibri}{\setmainfont{Calibri}}{%
  \IfFontExistsTF{Carlito}{\setmainfont{Carlito}}{\setmainfont{Liberation Sans}}}

\IfFontExistsTF{Arial}{\setsansfont{Arial}}{%
  \IfFontExistsTF{Liberation Sans}{\setsansfont{Liberation Sans}}{\setsansfont{DejaVu Sans}}}

\IfFontExistsTF{Courier New}{\setmonofont[Scale=0.9]{Courier New}}{%
  \IfFontExistsTF{Liberation Mono}{\setmonofont[Scale=0.9]{Liberation Mono}}{%
    \setmonofont[Scale=0.9]{DejaVu Sans Mono}}}
